\pagestyle{empty}
\tight
\begin{tblr}{width=\textwidth,colspec={|Q[c,m,1.9cm]|X[l,m]|},hlines,vlines,hspan=minimal,row{1}={bg=headblue},rowsep=1pt,colsep=3pt}
\SetCell{c,m}\hfil\logo\hfil & \textbf{POLITEKNIK NEGERI MALANG}\par\textbf{JURUSAN TEKNOLOGI INFORMASI}\par\textbf{PROGRAM STUDI: D4 TEKNIK INFORMATIKA} \\
\SetCell[c=2]{c} \textbf{RENCANA PEMBELAJARAN SEMESTER (RPS)} & \\
\end{tblr}
\begin{longtblr}{width=\textwidth,colspec={|Q[l,m,1.9cm]|X[0.85,l,m]|X[1.45,l,m]|X[0.75,l,m]|X[1.15,l,m]|},hlines,vlines,hspan=minimal,rowsep=1pt,colsep=2pt,row{1,3}={bg=headgray,font=\bfseries}}
MATA KULIAH & KODE & BOBOT (SKS)/jam & SEMESTER & TGL. PENYUSUNAN \\
Fisika & RTI251009 & 2 SKS/2 jam & 1 & 7 Agustus 2025 \\
OTORISASI & Dosen Pengembang RPS & Koordinator RMK & \SetCell[c=2]{l} Ka PRODI &  \\
 & Vipkas Al Hadid Firdaus, S.T., M.T.\par Sofyan Noor Arief, S.ST., M.Kom.\par Agung Nugroho Pramudhita, S.T., M.T. & Vipkas Al Hadid Firdaus, S.T., M.T. & \SetCell[c=2]{l}{\vspace{8mm}\par Dr. Ely Setyo Astuti, S.T., M.T.} &  \\
\SetCell[r=9]{l} Capaian Pembelajaran (CP) & \SetCell[c=4]{l,bg=headgray,font=\bfseries} Capaian Pembelajaran Lulusan Program Studi (CPL-Prodi) &  &  &  \\
 & CPL09 & \SetCell[c=3]{l} Mampu menjelaskan cara kerja sistem komputer dan menerapkan pendekatan komputasi untuk menyelesaikan masalah. &  &  \\
 & \SetCell[c=4]{l,bg=headgray,font=\bfseries} Capaian Pembelajaran mata kuliah (CPL-MK) &  &  &  \\
 & CPMK0902 & \SetCell[c=3]{l} Mahasiswa mampu menerapkan prinsip matematis dan struktur diskrit sebagai dasar analisis komputasional. &  &  \\
 & \SetCell[c=4]{l,bg=headgray,font=\bfseries} Kemampuan akhir tiap tahapan belajar (Sub-CPMK) &  &  &  \\
 & SCPMK0902-00901 & \SetCell[c=3]{l} Mahasiswa mampu menjelaskan hakikat fisika, peran fisika dalam teknologi informasi, dan metode ilmiah, serta melakukan pengukuran, konversi satuan, dan analisis ketidakpastian. &  &  \\
 & SCPMK0902-00902 & \SetCell[c=3]{l} Mahasiswa mampu menganalisis kinematika gerak lurus dan gerak dua dimensi serta menerapkan hukum Newton untuk menyelesaikan masalah dinamika partikel pada sistem mekanik. &  &  \\
 & SCPMK0902-00903 & \SetCell[c=3]{l} Mahasiswa mampu menganalisis hubungan usaha dan energi serta menerapkan konsep impuls, momentum, dan hukum kekekalan momentum pada permasalahan tumbukan. &  &  \\
 & SCPMK0902-00904 & \SetCell[c=3]{l} Mahasiswa mampu menganalisis gaya Coulomb, medan dan potensial listrik, kapasitansi, serta rangkaian listrik DC untuk permasalahan kelistrikan pada sistem elektronik. &  &  \\
\SetCell[r=2]{l} Penilaian & \SetCell[c=4]{l} *Nilai CPMK didapat dari agregasi nilai SUB-CPMK yang dibebankan &  &  &  \\
 & \SetCell[c=4]{l}{\tiny\begin{tblr}{width=\linewidth,colspec={|Q[c,m,0.1\linewidth]|Q[c,m,0.16\linewidth]|X[l,m]|X[l,m]|X[l,m]|X[l,m]|Q[c,m,0.08\linewidth]|},hlines,vlines,hspan=minimal,rowsep=1pt,colsep=0.7pt,row{1,2}={bg=headgray,font=\bfseries}}
Kode CPMK & Kode SCPMK & \SetCell[c=4]{c} Bobot per Bentuk Penilaian & & & & Total Bobot \\
 & & Tugas & Kuis & UTS & UAS & \\
\SetCell[r=4]{c} CPMK0902 & SCPMK0902-00901 & 5\%: konsep dasar, metode ilmiah, pengukuran & 10\%: Kuis 1 pekan 1--3 & 10\%: konsep dasar dan pengukuran & & 25\% \\
 & SCPMK0902-00902 & 6\%: kinematika dan dinamika partikel & & 20\%: kinematika dan hukum Newton & & 26\% \\
 & SCPMK0902-00903 & 3\%: usaha, energi, momentum & 10\%: Kuis 2 pekan 10--11 & & 8\%: usaha-energi dan tumbukan & 21\% \\
 & SCPMK0902-00904 & 6\%: kelistrikan dan rangkaian DC & & & 22\%: kelistrikan statis dan rangkaian DC & 28\% \\
\SetCell[c=6]{r}\textbf{Total Per Penilaian} & & & & & & \textbf{100\%} \\
\end{tblr}} &  &  &  \\
Diskripsi Singkat Mata Kuliah & \SetCell[c=4]{l} Mata kuliah Fisika memberikan pemahaman dasar mengenai konsep-konsep fisika klasik yang relevan dengan bidang teknik informatika. Topik meliputi pengukuran, mekanika, dinamika, energi, momentum, dan kelistrikan, dengan integrasi kasus-kasus dalam sistem komputer, Internet of Things, dan pemrosesan sinyal digital. &  &  &  \\
Bahan Kajian / Materi Pembelajaran & \SetCell[c=4]{l}\begin{enumerate}\item Hakikat dan ruang lingkup fisika\item Pengantar ilmu sains dan metode ilmiah\item Pengukuran dan satuan\item Kinematika gerak lurus dan dua dimensi\item Dinamika partikel\item Usaha dan energi\item Impuls dan momentum\item Medan listrik dan gaya Coulomb\item Potensial listrik\item Kapasitansi dan energi dalam medan listrik\item Arus listrik dan analisis rangkaian listrik DC\end{enumerate} &  &  &  \\
\SetCell[r=2]{l} Pustaka & \SetCell[c=4]{l} \textbf{Utama:}\begin{enumerate}\item Serway, R. A. \& Jewett, J. W. 2014. Physics for Scientists and Engineers with Modern Physics (9th Edition). Brooks/Cole -- Cengage Learning.\item Giancoli, D. C. 2005. Physics for Scientists and Engineers with Modern Physics (4th Edition). Pearson Education.\end{enumerate} &  &  &  \\
 & \SetCell[c=4]{l} \textbf{Pendukung:} Purwanto, S. \& Rusyidi, C. 2020. Fisika Dasar: Teori dan Aplikasi dalam Sains dan Teknik. Penerbit UPI Press; materi LMS dan lembar kerja. &  &  &  \\
Media Pembelajaran & \SetCell[c=2]{l} \textbf{Software:} LMS, aplikasi simulasi fisika, office suite. &  & \SetCell[c=2]{l} \textbf{Hardware:} Komputer/laptop, proyektor. &  \\
Nama Dosen Pengampu & \SetCell[c=4]{l}\begin{enumerate}\item Vipkas Al Hadid Firdaus, S.T., M.T.\item Sofyan Noor Arief, S.ST., M.Kom.\item Agung Nugroho Pramudhita, S.T., M.T.\end{enumerate} &  &  &  \\
Matakuliah Syarat & \SetCell[c=4]{l} - &  &  &  \\
\end{longtblr}
\begin{longtblr}{width=\textwidth,colspec={|Q[c,m,0.06\textwidth]|X[1.35,l,m]|X[1.08,l,m]|X[1.16,l,m]|X[0.7,l,m]|X[1.24,l,m]|X[1.06,l,m]|X[0.98,l,m]|Q[c,m,0.05\textwidth]|},hlines,vlines,hspan=minimal,rowhead=2,row{1,2}={bg=headgreen,font=\bfseries},rowsep=1pt,colsep=1.5pt}
Mgg.\par Ke & Kemampuan Akhir Yang Direncanakan (Sub-CP-MK) & Bahan kajian (Materi Pembelajaran) & Bentuk dan Metode Pembelajaran & Estimasi Waktu & Pengalaman Belajar Mahasiswa & Kriteria \& Bentuk Penilaian & Indikator Penilaian & Bobot Penilaian (\%) \\
(1) & (2) & (3) & (4) & (5) & (6) & (7) & (8) & (9) \\
1 & SCPMK0902-00901 & Hakikat dan ruang lingkup fisika; peran fisika dalam bidang TI. & Bentuk: Luring/Daring. Metode: Contextual Teaching and Learning (CTL) dan Problem Based Learning. & 2 x 2 x 100' tatap muka dan tugas terstruktur; 2 x 2 x 70' tugas mandiri. & Mahasiswa mempelajari hakikat dan ruang lingkup fisika, mengidentifikasi peran fisika dalam teknologi informasi, dan menyebutkan contoh penerapannya pada sistem teknologi modern. & Kriteria: ketepatan dan penguasaan. Bentuk: keaktifan diskusi kelompok dan ketepatan jawaban tugas. Mengacu rubrik RTM. & Menjelaskan pengertian fisika; mengaitkan fisika dan TI; menyebutkan contoh penerapan fisika. & 2 \\
2 & SCPMK0902-00901 & Metode ilmiah dan proses berpikir ilmiah. & Bentuk: Luring/Daring. Metode: Contextual Teaching and Learning (CTL) dan Problem Based Learning. & 2 x 2 x 100' tatap muka dan tugas terstruktur; 2 x 2 x 70' tugas mandiri. & Mahasiswa mempelajari langkah-langkah metode ilmiah, menganalisis masalah ilmiah sederhana secara sistematis, dan menerapkan metode ilmiah pada studi kasus berbasis eksperimen. & Kriteria: ketepatan dan penguasaan. Bentuk: keaktifan diskusi kelompok dan ketepatan jawaban tugas. Mengacu rubrik RTM. & Menjelaskan tahapan metode ilmiah dan menerapkannya pada kasus. & 2 \\
3 & SCPMK0902-00901 & Besaran dan satuan; alat ukur dan ketidakpastian pengukuran. & Bentuk: Luring/Daring. Metode: Contextual Teaching and Learning (CTL) dan Project Based Learning. & 1 x 2 x 100' tatap muka dan tugas terstruktur; 1 x 2 x 70' tugas mandiri. & Mahasiswa mempelajari konsep besaran dan satuan, menggunakan alat ukur untuk pengukuran yang tepat, serta menghitung ketidakpastian hasil pengukuran dan konversi satuan. & Kriteria: ketepatan dan penguasaan. Bentuk: keaktifan diskusi kelompok dan ketepatan jawaban tugas. Mengacu rubrik RTM. & Mengidentifikasi satuan; menghitung kesalahan; menyimpulkan hasil ukur. & 1 \\
4 & SCPMK0902-00901 & Kuis 1 materi pekan 1--3. & Ujian teori: pilihan ganda dan/atau esai, sifat closed book. & 1 x 2 x 100' ujian tertulis. & Mahasiswa mengerjakan soal kuis dengan materi pekan 1--3 secara mandiri. & Ujian tertulis; ketepatan jawaban dengan kunci jawaban. Mengacu rubrik RTM. & Ketepatan jawaban dengan kunci jawaban. & 10 \\
5 & SCPMK0902-00902 & Kinematika gerak lurus (GLB dan GLBB). & Bentuk: Luring/Daring. Metode: Contextual Teaching and Learning (CTL) dan Project Based Learning. & 2 x 2 x 100' tatap muka dan tugas terstruktur; 2 x 2 x 70' tugas mandiri. & Mahasiswa mempelajari konsep kecepatan dan percepatan gerak lurus, menggunakan persamaan gerak lurus, dan menganalisis grafik hubungan posisi, kecepatan, dan waktu. & Kriteria: ketepatan dan penguasaan. Bentuk: keaktifan diskusi kelompok dan ketepatan jawaban tugas. Mengacu rubrik RTM. & Menghitung dan menjelaskan GLB dan GLBB. & 1.5 \\
6 & SCPMK0902-00902 & Kinematika gerak dua dimensi. & Bentuk: Luring/Daring. Metode: Contextual Teaching and Learning (CTL) dan Project Based Learning. & 2 x 2 x 100' tatap muka dan tugas terstruktur; 2 x 2 x 70' tugas mandiri. & Mahasiswa menguraikan vektor posisi dan kecepatan dalam dua arah, menjelaskan gerak parabola, dan menghitung parameter gerak pada kasus nyata seperti peluru atau benda jatuh. & Kriteria: ketepatan dan penguasaan. Bentuk: keaktifan diskusi kelompok dan ketepatan jawaban tugas. Mengacu rubrik RTM. & Menghitung kecepatan, posisi, dan sudut lintasan gerak. & 1.5 \\
7 & SCPMK0902-00902 & Hukum Newton I, II, III dan gaya-gaya (normal, gesek, tegangan, sentripetal). & Bentuk: Luring/Daring. Metode: Contextual Teaching and Learning (CTL) dan Project Based Learning. & 2 x 2 x 100' tatap muka dan tugas terstruktur; 2 x 2 x 70' tugas mandiri. & Mahasiswa menjelaskan makna hukum Newton I, II, III, mengidentifikasi gaya-gaya yang bekerja pada benda, dan menerapkan hukum Newton untuk menentukan percepatan dan gaya total. & Kriteria: ketepatan dan penguasaan. Bentuk: keaktifan diskusi kelompok dan ketepatan jawaban tugas. Mengacu rubrik RTM. & Menentukan resultan gaya dan arah percepatan. & 1.5 \\
8 & SCPMK0902-00902 & Penerapan hukum Newton dalam sistem mekanik (rem otomatis, conveyor, lengan robotik). & Bentuk: Luring/Daring. Metode: Contextual Teaching and Learning (CTL) dan Project Based Learning. & 2 x 2 x 100' tatap muka dan tugas terstruktur; 2 x 2 x 70' tugas mandiri. & Mahasiswa menghubungkan gaya dan gerak dalam sistem teknologi, menganalisis peran hukum Newton pada rem otomatis dan lengan robot, serta menggunakannya dalam perancangan logika mekanik sederhana. & Kriteria: ketepatan dan penguasaan. Bentuk: keaktifan diskusi kelompok dan ketepatan jawaban tugas. Mengacu rubrik RTM. & Menghubungkan hukum Newton dengan fungsi sistem TI. & 1.5 \\
9 & SCPMK0902-00901, SCPMK0902-00902 & UTS materi pekan 1--8. & Ujian teori: pilihan ganda dan/atau esai, sifat closed book. & 1 x 2 x 100' ujian tertulis. & Mahasiswa mengerjakan soal ujian tengah semester dengan materi pekan 1--8 secara mandiri. & Ujian tertulis; ketepatan jawaban dengan kunci jawaban. Mengacu rubrik RTM. & Ketepatan jawaban dengan kunci jawaban. & 30 \\
10 & SCPMK0902-00903 & Usaha, energi kinetik, energi potensial, hukum kekekalan energi. & Bentuk: Luring/Daring. Metode: Contextual Teaching and Learning (CTL) dan Project Based Learning. & 2 x 2 x 100' tatap muka dan tugas terstruktur; 2 x 2 x 70' tugas mandiri. & Mahasiswa menjelaskan hubungan usaha dan perubahan energi, menghitung usaha oleh gaya konstan dan variabel, serta menggunakan hukum kekekalan energi untuk menganalisis sistem. & Kriteria: ketepatan dan penguasaan. Bentuk: keaktifan diskusi kelompok dan ketepatan jawaban tugas. Mengacu rubrik RTM. & Menganalisis konversi energi dalam sistem. & 1.5 \\
11 & SCPMK0902-00903 & Momentum, impuls, hukum kekekalan momentum, tumbukan elastis dan inelastis. & Bentuk: Luring/Daring. Metode: Contextual Teaching and Learning (CTL) dan Project Based Learning. & 2 x 2 x 100' tatap muka dan tugas terstruktur; 2 x 2 x 70' tugas mandiri. & Mahasiswa menjelaskan konsep impuls dan momentum linear, menganalisis tumbukan elastis dan inelastis, serta menerapkan hukum kekekalan momentum dalam simulasi tumbukan. & Kriteria: ketepatan dan penguasaan. Bentuk: keaktifan diskusi kelompok dan ketepatan jawaban tugas. Mengacu rubrik RTM. & Menghitung dan membandingkan hasil tumbukan. & 1.5 \\
12 & SCPMK0902-00903 & Kuis 2 materi pekan 10--11. & Ujian teori: pilihan ganda dan/atau esai, sifat closed book. & 1 x 2 x 100' ujian tertulis. & Mahasiswa mengerjakan soal kuis dengan materi pekan 10--11 secara mandiri. & Ujian tertulis; ketepatan jawaban dengan kunci jawaban. Mengacu rubrik RTM. & Ketepatan jawaban dengan kunci jawaban. & 10 \\
13 & SCPMK0902-00904 & Muatan listrik, hukum Coulomb, medan listrik. & Bentuk: Luring/Daring. Metode: Contextual Teaching and Learning (CTL) dan Project Based Learning. & 2 x 2 x 100' tatap muka dan tugas terstruktur; 2 x 2 x 70' tugas mandiri. & Mahasiswa menghitung gaya antara muatan menggunakan hukum Coulomb, menggambarkan garis medan listrik di sekitar muatan, dan menganalisis distribusi muatan pada sistem. & Kriteria: ketepatan dan penguasaan. Bentuk: keaktifan diskusi kelompok dan ketepatan jawaban tugas. Mengacu rubrik RTM. & Menghitung gaya Coulomb dan menggambar medan listrik. & 1.5 \\
14 & SCPMK0902-00904 & Potensial listrik dan hubungannya dengan medan listrik. & Bentuk: Luring/Daring. Metode: Contextual Teaching and Learning (CTL) dan Project Based Learning. & 2 x 2 x 100' tatap muka dan tugas terstruktur; 2 x 2 x 70' tugas mandiri. & Mahasiswa menjelaskan konsep energi potensial listrik dan beda potensial, menghubungkan medan listrik dengan potensial, dan menganalisis lintasan muatan dalam medan potensial. & Kriteria: ketepatan dan penguasaan. Bentuk: keaktifan diskusi kelompok dan ketepatan jawaban tugas. Mengacu rubrik RTM. & Menentukan energi potensial dan beda potensial. & 1.5 \\
15 & SCPMK0902-00904 & Kapasitor, kapasitansi, dan energi tersimpan. & Bentuk: Luring/Daring. Metode: Contextual Teaching and Learning (CTL) dan Project Based Learning. & 2 x 2 x 100' tatap muka dan tugas terstruktur; 2 x 2 x 70' tugas mandiri. & Mahasiswa menjelaskan konsep kapasitansi, menghitung kapasitansi dan energi yang tersimpan dalam kapasitor, serta menganalisis rangkaian kapasitor pada sistem elektronik. & Kriteria: ketepatan dan penguasaan. Bentuk: keaktifan diskusi kelompok dan ketepatan jawaban tugas. Mengacu rubrik RTM. & Menghitung kapasitansi dan energi dari konfigurasi kapasitor. & 1.5 \\
16 & SCPMK0902-00904 & Rangkaian listrik DC sederhana (hukum Ohm dan hukum Kirchhoff). & Bentuk: Luring/Daring. Metode: Contextual Teaching and Learning (CTL) dan Project Based Learning. & 2 x 2 x 100' tatap muka dan tugas terstruktur; 2 x 2 x 70' tugas mandiri. & Mahasiswa menjelaskan hukum Ohm dan hukum Kirchhoff, menganalisis rangkaian seri dan paralel, serta menghitung arus, tegangan, daya, dan efisiensi rangkaian. & Kriteria: ketepatan dan penguasaan. Bentuk: keaktifan diskusi kelompok dan ketepatan jawaban tugas. Mengacu rubrik RTM. & Menghitung arus, tegangan, daya, dan efisiensi. & 1.5 \\
17 & SCPMK0902-00903, SCPMK0902-00904 & UAS materi pekan 1--16. & Ujian teori: pilihan ganda dan/atau esai, sifat closed book. & 1 x 2 x 100' ujian tertulis. & Mahasiswa mengerjakan soal ujian akhir semester dengan materi pekan 1--16 secara mandiri. & Ujian tertulis; ketepatan jawaban dengan kunci jawaban. Mengacu rubrik RTM. & Ketepatan jawaban dengan kunci jawaban. & 30 \\
\end{longtblr}
